\SourceLeaf{111}%
\begin{center}$+3$ pertinet\end{center}
{\small
\[
\begin{array}{c|l}
\substack{\text{ad}\\\text{complexum}}&\text{si}\\
A&\begin{aligned}b\equiv0\ (\mathrm{mod.}\ 12);\\\text{vel simul }a\equiv0\ (\mathrm{mod.}\ 3),\ b\equiv2\ (\mathrm{mod.}\ 4)\end{aligned}\\
B&b\equiv8a\text{ vel }10a\ (\mathrm{mod.}\ 12)\\
C&\begin{aligned}b\equiv6a\ (\mathrm{mod.}\ 12);\\\text{vel simul }a\equiv0\ (\mathrm{mod.}\ 3),\ b\equiv0\ (\mathrm{mod.}\ 4)\end{aligned}\\
D&b\equiv2a\text{ vel }4a\ (\mathrm{mod.}\ 12)
\end{array}
\]
}

Perinde criteria pro $\pm6$ petuntur e combinatione criteriorum pro $\mp2$ et $-3$; scilicet
\begin{center}$+6$ pertinet\end{center}
{\small
\[
\begin{array}{c|l}
\substack{\text{ad}\\\text{complexum}}&\text{si}\\
A&\begin{aligned}b\equiv0,\quad2a,22a\ (\mathrm{mod.}\ 24);\\\text{vel simul }a\equiv0\ (\mathrm{mod.}\ 3),\ b\equiv4a\ (\mathrm{mod.}\ 8)\end{aligned}\\
B&\begin{aligned}b\equiv4a,\quad6a,\quad8a\ (\mathrm{mod.}\ 24);\\\text{vel simul }a\equiv0\ (\mathrm{mod.}\ 3),\ b\equiv2a\ (\mathrm{mod.}\ 8)\end{aligned}\\
C&\begin{aligned}b\equiv10a,12a,14a\ (\mathrm{mod.}\ 24);\\\text{vel simul }a\equiv0\ (\mathrm{mod.}\ 3),\ b\equiv0\ (\mathrm{mod.}\ 8)\end{aligned}\\
D&\begin{aligned}b\equiv16a,18a,20a\ (\mathrm{mod.}\ 24);\\\text{vel simul }a\equiv0\ (\mathrm{mod.}\ 3),\ b\equiv6a\ (\mathrm{mod.}\ 8)\end{aligned}
\end{array}
\]
}

\begin{center}$-6$ vero\end{center}
{\small
\[
\begin{array}{c|l}
\substack{\text{ad}\\\text{complexum}}&\text{si}\\
A&\begin{aligned}b\equiv0,\quad10a,14a\ (\mathrm{mod.}\ 24);\\\text{vel simul }a\equiv0\ (\mathrm{mod.}\ 3),\ b\equiv4a\ (\mathrm{mod.}\ 8)\end{aligned}\\
B&\begin{aligned}b\equiv4a,\quad8a,18a\ (\mathrm{mod.}\ 24);\\\text{vel simul }a\equiv0\ (\mathrm{mod.}\ 3),\ b\equiv6a\ (\mathrm{mod.}\ 8)\end{aligned}\\
C&\begin{aligned}b\equiv2a,12a,22a\ (\mathrm{mod.}\ 24);\\\text{vel simul }a\equiv0\ (\mathrm{mod.}\ 3),\ b\equiv0\ (\mathrm{mod.}\ 8)\end{aligned}\\
D&\begin{aligned}b\equiv6a,16a,20a\ (\mathrm{mod.}\ 24);\\\text{vel simul }a\equiv0\ (\mathrm{mod.}\ 3),\ b\equiv2a\ (\mathrm{mod.}\ 8)\end{aligned}
\end{array}
\]
}

Simili modo criteria pro numero $+21$ concinnabuntur e criteriis pro $-3$ et $-7$; criteria pro $-105$ e criteriis pro $-1$, $-3$, $+5$, $-7$, etc.

\Article{30}
Amplissimam itaque messem theorematum specialium aperit inductio, theoremati pro numero $2$ affinium: sed desideratur vinculum commune, desiderantur demonstrationes rigorosae, quum methodus, per quam in commentatione prima numerum $2$ absolvimus, ulteriorem applicationem non patiatur. Non desunt quidem methodi diversae, per quas demonstrationibus pro casibus particularibus potiri liceret, iis potissimum, qui distributionem residuorum quadraticorum inter complexus $A,C$ spectant, quibus tamen non immoramur, quum theoria genera%
\SourceLeaf{112}%
lis \emph{omnes} casus complectens in votis esse debeat. Cui rei quum inde ab anno $1805$ meditationes nostras dicare coepissemus, mox certiores facti sumus, fontem genuinum theoriae generalis in campo arithmeticae promoto quaerendum esse, uti iam in art. 1 addigitavimus.

Quemadmodum scilicet arithmetica sublimior in quaestionibus hactenus pertractatis inter solos numeros integros reales versatur, ita theoremata circa residua biquadratica tunc tantum in summa simplicitate ac genuina venustate resplendent, quando campus arithmeticae ad quantitates \emph{imaginarias} extenditur, ita ut absque restrictione ipsius obiectum constituant numeri formae $a+bi$, denotantibus $i$ pro more quantitatem imaginariam $\surd-1$, atque $a,b$ indefinite omnes numeros reales integros inter $-\infty$ et $+\infty$. Tales numeros vocabimus \emph{numeros integros complexos}, ita quidem, ut reales complexis non opponantur, sed tamquam species sub his contineri censeantur. Commentatio praesens tum doctrinam elementarem de numeris complexis, tum prima initia theoriae residuorum biquadraticorum sistet, quam ab omni parte perfectam reddere in continuatione subsequente suscipiemus\begingroup\renewcommand{\thefootnote}{*)}\footnote{Obiter saltem hic adhuc monere convenit, campum ita definitum imprimis theoriae residuorum biquadraticorum accommodatum esse. Theoria residuorum cubicorum simili modo superstruenda est considerationi numerorum formae $a+bh$, ubi $h$ est radix imaginaria aequationis $h^3-1=0$, puta $h=-\frac12+\surd\frac34\pdot i$; et perinde theoria residuorum potestatum altiorum introductionem aliarum quantitatum imaginariarum postulabit.}\endgroup.

\Article{31}
Ante omnia quasdam denominationes praemittimus, per quarum introductionem brevitati et perspicuitati consuletur.

Campus numerorum complexorum $a+bi$ continet
\begin{list}{}{\setlength{\leftmargin}{1.6em}\setlength{\labelwidth}{1.4em}\setlength{\labelsep}{.2em}}
\item[I.] numeros reales, ubi $b=0$, et, inter hos, pro indole ipsius $a$
\begin{list}{}{\setlength{\leftmargin}{1.7em}\setlength{\labelwidth}{1.5em}\setlength{\labelsep}{.2em}}
\item[1)] cifram
\item[2)] numeros positivos
\item[3)] numeros negativos
\end{list}
\item[II.] numeros imaginarios, ubi $b$ cifrae inaequalis. Hic iterum distinguuntur
\begin{list}{}{\setlength{\leftmargin}{1.7em}\setlength{\labelwidth}{1.5em}\setlength{\labelsep}{.2em}}
\item[1)] numeri imaginarii absque parte reali, i. e. ubi $a=0$
\item[2)] numeri imaginarii cum parte reali, ubi neque $b$ neque $a=0$.
\end{list}
\end{list}
Priores si placet numeri imaginarii puri, posteriores numeri imaginarii mixti vocari possunt.

\SourceLeaf{113}%
Unitatibus in hac doctrina utimur quaternis, $+1$, $-1$, $+i$, $-i$, quae simpliciter positiva, negativa, positiva imaginaria, negativa imaginaria audient.

Producta terna cuiuslibet numeri complexi per $-1$, $+i$, $-i$ illius \emph{socios} vel \emph{numeros illi associatos} appellabimus. Excepta itaque cifra (quae sibi ipsa associata est), semper quaterni numeri \emph{inaequales} associati sunt.

Contra numero complexo \emph{coniunctum} vocamus eum, qui per permutationem ipsius $i$ cum $-i$ inde oritur. Inter numeros imaginarios itaque bini \emph{inaequales} semper coniuncti sunt, dum numeri reales sibi ipsi sunt coniuncti, siquidem denominationem ad hos extendere placet.

Productum numeri complexi per numerum ipsi coniunctum utriusque \emph{normam} vocamus. Pro norma itaque numeri realis, ipsius quadratum habendum est.

Generaliter octonos numeros nexos habemus, puta
\[
\begin{array}{r|r}
a+bi&a-bi\\
-b+ai&-b-ai\\
-a-bi&-a+bi\\
b-ai&b+ai
\end{array}
\]
ubi duas quaterniones numerorum associatorum, quatuor biniones coniunctorum conspicimus, omniumque norma communis est $aa+bb$. Sed octo numeri ad quatuor inaequales reducuntur, quoties vel $a=\pm b$, vel alteruter numerorum $a,b=0$.

E definitionibus allatis protinus demanant sequentia:

Producto duorum numerorum complexorum coniunctum est productum e numeris, qui illis coniuncti sunt.

Idem valet de producto e pluribus factoribus, nec non de quotientibus.

Norma producti e duobus numeris complexis aequalis est producto ex horum normis.

Hoc quoque theorema extenditur ad producta e quotcunque factoribus et ad quotientes.

Cuiusvis numeri complexi (excipiendo cifram, quod plerumque abhinc tacite subintelligemus) norma est numerus \emph{positivus}.

Ceterum nihil obstat, quominus definitiones nostrae ad valores fractos vel adeo irrationales ipsorum $a,b$ extendantur; sed $a+bi$ tunc tantum numerus complexus integer audiet, quando \emph{uterque} $a,b$ est integer; atque tunc tantum rationalis, quando \emph{uterque} $a,b$ rationalis est.

\SourceLeaf{114}%
\Article{32}
Algorithmus operationum arithmeticarum circa numeros complexos vulgo notus est: divisio, per introductionem normae, ad multiplicationem reducitur, quum habeatur
\[
\frac{a+bi}{c+di}=(a+bi)\pdot\frac{c-di}{cc+dd}=\frac{ac+bd}{cc+dd}+\frac{bc-ad}{cc+dd}\pdot i
\]

Extractio radicis quadratae perficitur adiumento formulae
\[
\surd(a+bi)=\pm\left(\surd\frac{\surd(aa+bb)+a}{2}+i\surd\frac{\surd(aa+bb)-a}{2}\right)
\]
si $b$ est numerus positivus, vel huius
\[
\surd(a+bi)=\pm\left(\surd\frac{\surd(aa+bb)+a}{2}-i\surd\frac{\surd(aa+bb)-a}{2}\right)
\]
si $b$ est numerus negativus. Usui transformationis quantitatis complexae $a+bi$ in $r(\cos\varphi+i\sin\varphi)$ ad calculos facilitandos, non opus est hic immorari.

\Article{33}
Numerum integrum complexum, qui in factores duos ab unitatibus diversos\begingroup\renewcommand{\thefootnote}{*)}\footnote{sive, quod idem est, tales, quorum normae unitate sint maiores.}\endgroup\ resolvi potest, vocamus numerum complexum compositum; contra numerus primus complexus dicetur, qui talem resolutionem in factores non admittit. Hinc statim patet, quemvis numerum compositum realem etiam esse compositum complexum. At numerus primus realis poterit esse numerus complexus compositus, et quidem hoc valebit de numero $2$ atque de omnibus numeris primis realibus positivis formae $4n+1$ (excepto numero $1$), quippe quos in bina quadrata positiva decomponi posse constat; puta, fit $2=(1+i)(1-i)$, $5=(1+2i)(1-2i)$, $13=(3+2i)(3-2i)$, $17=(1+4i)(1-4i)$ etc.

Contra numeri primi reales positivi formae $4n+3$ semper sunt numeri primi complexi. Si enim talis numerus $q$ esset $=(a+bi)(\alpha+\beta i)$, foret etiam $q=(a-bi)(\alpha-\beta i)$, adeoque $qq=(aa+bb)(\alpha\alpha+\beta\beta)$: at $qq$ unico tantum modo in factores positivos unitate maiores resolvi potest, puta in $q\times q$, unde esse debere $q=aa+bb=\alpha\alpha+\beta\beta$, Q. E. A.; quum summa duorum quadratorum nequeat esse formae $4n+3$.

\SourceLeaf{115}%
Numeri reales negativi manifesto easdem denominationes servant, quas positivi, idemque valet de numeris imaginariis puris.

Superest itaque, ut inter numeros imaginarios mixtos, compositos a primis dignoscere doceamus, quod fit per sequens

\textsc{Theorema.}\quad\emph{Quivis numerus integer imaginarius mixtus $a+bi$ est vel numerus primus complexus, vel numerus compositus, prout ipsius norma est vel numerus primus realis, vel numerus compositus.}

\emph{Dem.} I. Quoniam numeri complexi compositi norma semper est numerus compositus, patet, numerum complexum, cuius norma sit numerus primus realis, necessario esse debere numerum primum complexum. Q. E. P.

II. Si vero norma $aa+bb$ est numerus compositus, sit $p$ numerus primus positivus realis illam metiens. Duo iam casus distinguendi sunt.

1) Si $p$ est formae $4n+3$, constat, $aa+bb$ per $p$ divisibilem esse non posse, nisi $p$ simul metiatur ipsos $a,b$, unde $a+bi$ erit numerus compositus.

2) Si $p$ non est formae $4n+3$, certo in duo quadrata decomponi poterit: statuemus itaque $p=\alpha\alpha+\beta\beta$. Quum fiat
\[
(a\alpha+b\beta)(a\alpha-b\beta)=aa(\alpha\alpha+\beta\beta)-\beta\beta(aa+bb)
\]
adeoque per $p$ divisibilis, $p$ certo alterutrum factorem $a\alpha+b\beta$, $a\alpha-b\beta$ metietur, et quum insuper fiat
\[
(a\alpha+b\beta)^2+(b\alpha-a\beta)^2=(a\alpha-b\beta)^2+(b\alpha+a\beta)^2=(aa+bb)(\alpha\alpha+\beta\beta)
\]
adeoque per $pp$ divisibilis, patet, in casu priori etiam $b\alpha-a\beta$, in posteriori $b\alpha+a\beta$ per $p$ divisibilem esse debere. Quare in casu priori
\[
\frac{a+bi}{\alpha+\beta i}=\frac{a\alpha+b\beta}{p}+\frac{b\alpha-a\beta}{p}\pdot i
\]
erit numerus integer complexus, in posteriori autem
\[
\frac{a+bi}{\alpha-\beta i}=\frac{a\alpha-b\beta}{p}+\frac{b\alpha+a\beta}{p}\pdot i
\]
integer erit. Quum itaque numerus propositus vel per $\alpha+\beta i$ vel per $\alpha-\beta i$ divisibilis sit, quotientisque norma $=\frac{aa+bb}{p}$ per hyp. ab unitate diversa fiat, patet, $a+bi$ in utroque casu esse numerum complexum compositum. Q. E. S.

\SourceLeaf{116}%
\Article{34}
Totum itaque ambitum numerorum primorum complexorum exhauriunt quatuor species sequentes:

\begin{list}{}{\setlength{\leftmargin}{1.7em}\setlength{\labelwidth}{1.5em}\setlength{\labelsep}{.2em}}
\item[1)] quatuor unitates, $1,+i,-1,-i$, quas tamen, dum de numeris primis agemus, plerumque tacite subintelligemus exclusas.
\item[2)] numerus $1+i$ cum tribus sociis $-1+i$, $-1-i$, $1-i$.
\item[3)] numeri primi reales positivi formae $4n+3$ cum ternis sociis.
\item[4)] numeri complexi, quorum normae sunt numeri primi reales formae $4n+1$ unitate maiores, et quidem cuivis normae tali datae semper octoni numeri primi complexi et non plures respondebunt, quum talis norma unico tantum modo in bina quadrata decomponi possit.
\end{list}

\Article{35}
Quemadmodum numeri integri reales in pares et impares distribuuntur, atque illi iterum in pariter pares et impariter pares, ita inter numeros complexos distinctio aeque essentialis se offert: sunt scilicet

vel per $1+i$ non divisibiles, puta numeri $a+bi$, ubi alter numerorum $a,b$ est impar, alter par;

vel per $1+i$ neque vero per $2$ divisibiles, quoties uterque $a,b$ est impar;

vel per $2$ divisibiles, quoties uterque $a,b$ est par.

Numeri primae classis commode dici possunt \emph{numeri complexi impares}, secundae \emph{semipares}, tertiae \emph{pares}.

Productum e pluribus factoribus complexis semper impar erit, quoties omnes factores sunt impares; semipar, quoties unus factor est semipar, reliqui impares; par autem, quoties inter factores vel saltem duo semipares inveniuntur, vel saltem unus par.

Norma cuiusvis numeri complexi imparis est formae $4n+1$; norma numeri semiparis est formae $8n+2$; denique norma numeri paris est productum numeri formae $4n+1$ in numerum $4$ vel altiorem binarii potestatem.

\Article{36}
Quum nexus inter quaternos numeros complexos socios analogus sit nexui inter binos numeros reales oppositos (i. e. absolute aequales signisque oppositis affectos), atque ex his vulgo positivus tamquam primarius merito considerari soleat:
\SourceLeaf{117}%
quaestio oritur, num similis distinctio inter quaternos numeros complexos socios stabiliri possit, et pro utili haberi debeat. Ad quam decidendam perpendere oportet, principium distinctionis ita comparatum esse debere, ut productum duorum numerorum, qui inter socios suos pro primariis valent, semper fiat numerus primarius inter socios suos. At mox certiores fimus, tale principium omnino non dari, nisi distinctio ad numeros integros restringatur: quinadeo distinctio \emph{utilis} ad numeros impares limitanda erit. Pro his vero finis propositus duplici modo attingi potest. Scilicet

I. Productum duorum numerorum $a+bi,a'+b'i$ ita comparatorum, ut $a,a'$ sint formae $4n+1$, atque $b,b'$ pares, eadem proprietate gaudebit, ut pars realis fiat $\equiv1\ (\mathrm{mod.}\ 4)$, atque pars imaginaria par. Et facile perspicitur, inter quaternos numeros impares associatos unum solum sub illa forma contentum esse.

II. Si numerus $a+bi$ ita comparatus est, ut $a-1$ et $b$ vel simul pariter pares sint, vel simul impariter pares, eius productum per numerum complexum eiusdem formae eadem forma gaudebit, facileque perspicitur, e quaternis numeris imparibus associatis unum solum sub hac forma contineri.

Ex his duobus principiis aeque fere idoneis posterius adoptabimus, scilicet inter quaternos numeros complexos impares associatos eum pro primario habebimus, qui secundum modulum $2+2i$ unitati positivae fit congruus: hoc pacto plura insignia theoremata maiori concinnitate enunciare licebit. Ita e. g. sunt numeri primi complexi primarii $-1+2i$, $-1-2i$, $+3+2i$, $+3-2i$, $+1+4i$, $+1-4i$ etc., nec non reales $-3$, $-7$, $-11$, $-19$ etc. manifesto semper signo negativo afficiendi. Numero complexo impari primario coniunctus quoque primarius erit.

Pro numeris semiparibus et paribus in genere similis distinctio nimis arbitraria parumque utilis foret. E numeris primis associatis $1+i$, $1-i$, $-1+i$, $-1-i$ unum quidem prae reliquis pro primario eligere possumus, sed ad compositos talem distinctionem non extendemus.

\Article{37}
Si inter factores numeri complexi compositi inveniuntur tales, qui ipsi sunt compositi, atque hi iterum in factores suos resolvuntur, manifesto tandem ad factores primos delabimur, i. e. quivis numerus compositus in factores primos resolubilis est. Inter quos si qui non primarii reperiuntur, singulorum loco substitua%
\SourceLeaf{118}%
tur productum primarii associati per $i,-1$ vel $-i$. Hoc pacto patet, quemvis numerum complexum compositum $M$ reduci posse ad formam
\[
M=i^\mu A^\alpha B^\beta C^\gamma\dotsfour
\]
ita ut $A,B,C$ etc. sint numeri primi complexi primarii inaequales, atque $\mu=0,1,2$ vel $3$. Circa hanc resolutionem theorema se offert, unico tantum modo eam fieri posse, quod theorema obiter quidem consideratum per se manifestum videri posset, sed utique demonstratione eget. Ad quam sternit viam sequens

\textsc{Theorema.}\quad\emph{Productum $M=A^\alpha B^\beta C^\gamma\dotsfour$, denotantibus $A,B,C$ etc. numeros primos complexos primarios diversos, divisibile esse nequit per ullum numerum primum complexum primarium, qui inter $A,B,C$ etc. non reperitur.}

\emph{Dem.} Sit $P$ numerus primus complexus primarius inter $A,B,C$ etc. non contentus, sintque $p,a,b,c$ etc. normae numerorum $P,A,B,C$ etc. Hinc facile colligitur, normam numeri $M$ fore $=a^\alpha b^\beta c^\gamma$ etc., unde hic numerus, si $M$ per $P$ divisibilis esset, per $p$ divisibilis esse deberet. Quum singulae normae sint vel numeri primi reales (e serie $2,5,13,17$ etc.), vel numerorum primorum realium quadrata (e serie $9,49,121$ etc.), sponte patet, illud evenire non posse, nisi $p$ cum aliqua norma $a,b,c$ etc. identica fiat: supponemus itaque $p=a$. At quum $P,A$ per hyp. sint numeri primi complexi primarii non identici, facile perspicitur, haec simul consistere non posse, nisi $P,A$ sint numeri complexi imaginarii coniuncti, et proin $p=a$ numerus primus realis impar, (non quadratum numeri primi): supponemus itaque $A=k+li$, $P=k-li$. Hinc (extendendo notionem et signum congruentiae ad numeros integros complexos) erit $A\equiv2k\ (\mathrm{mod.}\ P)$, unde facile colligitur
\[
M\equiv2^\alpha k^\alpha B^\beta C^\gamma\dotsfour\ (\mathrm{mod.}\ P)
\]
Quapropter dum $M$ per $P$ divisibilis supponitur, erit etiam
\[
2^\alpha k^\alpha B^\beta C^\gamma\dotsfour
\]
per $P$ divisibilis, adeoque norma huius numeri, quae fit
\[
=2^{2\alpha}k^{2\alpha}b^\beta c^\gamma\dotsfour
\]
divisibilis per $p$. At quum $2$ et $k$ per $p$ certo non sint divisibiles, hinc sequi%
\SourceLeaf{119}%
tur, $p$ cum aliquo numerorum $b,c$ etc. identicum esse debere: sit e. g. $p=b$. Hinc vero concludimus, esse vel $B=k+li$, vel $B=k-li$, i. e. vel $B=A$, vel $B=P$, utrumque contra hyp.

Ex hoc theoremate alterum, quod resolutio in factores primos unico tantum modo perfici potest, facillime derivatur, et quidem per ratiocinia iis, quibus in \emph{Disquisitionibus Arithmeticis} pro numeris realibus usi sumus (art. 16), prorsus analoga: quapropter illis hic immorari superfluum foret.

\Article{38}
Progredimur iam ad congruentiam numerorum secundum modulos complexos. Sed in limine huius disquisitionis convenit indicare, quomodo ditio quantitatum complexarum intuitui subiici possit.

Sicuti omnis quantitas realis per partem rectae utrinque infinitae ab initio arbitrario sumendam, et secundum segmentum arbitrarium pro unitate acceptum aestimandam exprimi, adeoque per punctum alterum repraesentari potest, ita ut puncta ab altera initii plaga quantitates positivas, ab altera negativas repraesentent: ita quaevis quantitas complexa repraesentari poterit per aliquod punctum in plano infinito, in quo recta determinata ad quantitates reales refertur, scilicet quantitas complexa $x+iy$ per punctum, cuius abscissa $=x$, ordinata (ab altera lineae abscissarum plaga positive, ab altera negative sumta) $=y$. Hoc pacto dici potest, quamlibet quantitatem complexam mensurare inaequalitatem inter situm puncti ad quod refertur atque situm puncti initialis, denotante unitate positiva deflexum arbitrarium determinatum versus directionem arbitrariam determinatam; unitate negativa deflexum aeque magnum versus directionem oppositam; denique unitatibus imaginariis deflexus aeque magnos versus duas directiones laterales normales.

Hoc modo metaphysica quantitatum, quas imaginarias dicimus, insigniter illustratur. Si punctum initiale per $(0)$ denotatur, atque duae quantitates complexae $m,m'$ ad puncta $M,M'$ referuntur, quorum situm relative ad $(0)$ exprimunt, differentia $m-m'$ nihil aliud erit nisi situs puncti $M$ relative ad punctum $M'$: contra, producto $mm'$ repraesentante situm puncti $N$ relative ad $(0)$, facile perspicies, hunc situm perinde determinari per situm puncti $M$ ad $(0)$, ut situs puncti $M'$ determinatur per situm puncti cui respondet unitas positiva, ita ut haud inepte dicas, situs punctorum respondentium quantitatibus complexis $mm'$
\SourceLeaf{120}%
$m,m',1$ formare \emph{proportionem}. Sed uberiorem huius rei tractationem ad aliam occasionem nobis reservamus. Difficultates, quibus theoria quantitatum imaginariarum involuta putatur, ad magnam partem a denominationibus parum idoneis originem traxerunt (quum adeo quidam usi sint nomine absono quantitatum impossibilium), Si, a conceptibus, quos offerunt varietates duarum dimensionum, (quales in maxima puritate conspiciuntur in intuitionibus spatii) profecti, quantitates positivas directas, negativas inversas, imaginarias laterales nuncupavissemus, pro tricis simplicitas, pro caligine claritas successisset.

\Article{39}
Quae in art. praec. prolata sunt, ad quantitates complexas continuas referuntur: in arithmetica, quae tantummodo circa numeros integros versatur, schema numerorum complexorum erit systema punctorum aequidistantium et in rectis aequidistantibus ita dispositorum, ut planum infinitum in infinite multa quadrata aequalia dispertiant. Omnes numeri per numerum complexum datum $a+bi=m$ divisibiles item infinite multa quadrata formabunt, quorum latera $=\surd(aa+bb)$ sive areae $=aa+bb$; quadrata posteriora ad priora inclinata erunt, quoties quidem neuter numerorum $a,b$ est $=0$. Cuivis numero per modulum $m$ non divisibili respondebit punctum vel intra tale quadratum situm vel in latere duobus quadratis contiguo; posterior tamen casus locum habere nequit, nisi $a,b$ divisorem communem habent: porro patet, numeros secundum modulum $m$ congruos in quadratis suis locos congruentes occupare. Hinc facile concluditur, si colligantur omnes numeri intra quadratum determinatum siti, nec non omnes qui forte in duobus eius lateribus non oppositis iaceant, denique his adscribatur numerus per $m$ divisibilis, haberi systema completum residuorum incongruorum secundum modulum $m$, i. e. quemvis integrum alicui ex illis et quidem unico tantum congruum esse debere. Nec difficile foret ostendere, horum residuorum multitudinem aequalem esse moduli normae, puta $=aa+bb$. Sed consultum videtur, hoc gravissimum theorema alio modo pure arithmetico demonstrare.

\Article{40}
\textsc{Theorema.}\quad\emph{Secundum modulum complexum datum $m=a+bi$, cuius norma $aa+bb=p$, et pro quo $a,b$ sunt numeri inter se primi, quilibet integer complexus congruus erit alicui residuo e serie $0,1,2,3\dotsfour p-1$, et non pluribus.}

